Beim Rutschen bremst die Gleitreibungskraft $\mu_\mathrm{G}\,m\,g$ den Velofahrer gleichmässig mit $a = \mu_\mathrm{G}\,g$. Seine mittlere Geschwindigkeit ist $v/2$, für die Bremsspur gilt also
\begin{align}
 s &= \frac{v}{2}\cdot t = \frac{v}{2}\cdot\frac{v}{a} = \frac{v^2}{2\,\mu_\mathrm{G}\,g}
\end{align}
Daraus folgt
\begin{align}
 \mu_\mathrm{G} &= \muGF \\
   &= \frac{\left(\vVe\right)^2}{2\cdot\gE\cdot\sBr} \\
   &= \muG \approx \result{\muGP{0}{2}}
\end{align}

Ohne Rutschen bremst die Haftreibung. Für Pneu auf trockener Asphaltstrasse gibt FoTaBe (S.~189) $\mu_\mathrm{H} = \muHO\pm\dmuO$; optimale Bedingungen bedeuten den grössten Wert $\mu_\mathrm{H,max} = \muHm$. Mit derselben Überlegung wie oben ist der Bremsweg
\begin{align}
 s &= \sOpF \\
   &= \frac{\left(\vVe\right)^2}{2\cdot\muHm\cdot\gE} \\
   &= \sOp \approx \result{\sOpP{0}{2}}
\end{align}
Mit dem typischen Wert $\mu_\mathrm{H} = \muHO$ wären es $\sTyP{0}{2}$. Die Reaktionszeit ist dabei nicht berücksichtigt.
